\section{Measure-valued weighted Caro--Wei}\label{app:upper}

We prove Lemma~\ref{lem:measure-cw} without requiring a jointly measurable family of independent
random clocks on an uncountable domain.

\begin{proof}[Proof of Lemma~\ref{lem:measure-cw}]
The claim is invariant to scaling $\mu$, so suppose first that $\mu$ is a probability measure.
Draw $Z_1,\ldots,Z_k$ independently from $\mu$. Construct a finite undirected graph on indices
$[k]$ by joining $i\ne j$ when $Z_i=Z_j$ or $Z_i$ and $Z_j$ are adjacent in $H$. Any independent
set of indices maps injectively to an independent set in $H$, so this sampled graph has independence
number at most $\alpha(H)$.

The ordinary Caro--Wei inequality applied conditionally on the sampled graph gives
\[
 \sum_{i=1}^k\frac{1}{1+d_i}\le\alpha(H),
\]
where $d_i$ is the degree of index $i$. Given $Z_i=x$, the closed-neighborhood count
$1+d_i$ has distribution $1+Y$ for
$Y\sim\operatorname{Bin}(k-1,p_x)$ and $p_x=\mu(N_H[x])$. Direct integration yields
\[
 \E\frac{1}{1+Y}
 =\int_0^1(1-p_x+p_xt)^{k-1}dt
 =\frac{1-(1-p_x)^k}{kp_x}.
\]
Taking expectation, summing the $k$ equal terms, and using measurability gives
\[
 \int f_k(p_x)\,d\mu(x)\le\alpha(H),\qquad
 f_k(p)=\begin{cases}[1-(1-p)^k]/p,&p>0,\\ k,&p=0.\end{cases}
\]
The functions $f_k$ increase pointwise to $1/p$ for $p>0$ and to infinity for $p=0$. Monotone
convergence therefore shows at once that the set $\{x:p_x=0\}$ is $\mu$-null and that the limiting
integral is at most $\alpha(H)$. For a finite nonzero measure, normalize by its total mass. The zero
measure is trivial.
\end{proof}

\begin{remark}
For a finite atomic measure, an equivalent proof assigns an independent exponential clock of rate
$\mu(x)$ to each positive-mass vertex. The set of vertices whose clock is smallest in its closed
neighborhood is independent and has expected size equal to the left side of
\eqref{eq:measure-cw}.
\end{remark}

\section{Directed lower bound}\label{app:directed}

We give a complete proof of Theorem~\ref{thm:directed-lower}. It suffices to consider deterministic
learners because independent learner randomness can be conditioned upon throughout.

\subsection{Posterior variance and harmonic weights}

Fix $q,m$ and set $s=\lceil m/q\rceil$, $n=qs$. Then $m\le n<m+q\le2m$ when $m\ge q$.
Use the marginal family \eqref{eq:paired-marginal} with $\eta=1/2$ and put the uniform prior on its
$n$ signs. Augment the observation by revealing the sign of a pair whenever either endpoint is
sampled. This can only decrease the minimax risk.

The total probability of every pair is $1/n$, independent of its sign. Thus the sequence of pair
indices in the sample is i.i.d. uniform on $[n]$, independent of all signs. Conditional on these
indices and the revealed signs, signs of unobserved pairs remain independent Rademachers.

For chain $a$, index suffixes by their number $r\in[s]$ of pairs. Let $U_{a,r}$ be the number of
unobserved pairs in that suffix. Equation~\eqref{eq:suffix-target} shows that its random part at
the first vertex is
\[
 -\frac{\eta}{2r}\sum_{j\in M_{a,r}}\sigma_j,
\]
where $M_{a,r}$ is the missing set. Its posterior variance is $\eta^2U_{a,r}/(4r^2)$. Also
$D_\sigma(u_{a,k})\ge(1-\eta)/(2n)$. Restricting the population loss to these first vertices proves
the conditional variance lower bound \eqref{eq:posterior-v}.

Let
\[
 w_t=\sum_{r=t}^s\frac1{r^2},\qquad
 Z=\sum_{a=1}^q\sum_{t=1}^s w_t I_{a,t},
\]
where $I_{a,t}$ indicates that the pair $t$ places from the end of chain $a$ is missing. Exchanging
the order of summation gives
\begin{equation}
 \sum_{a=1}^q\sum_{r=1}^s\frac{U_{a,r}}{r^2}=Z,
 \qquad
 \sum_{t=1}^s w_t=\sum_{r=1}^s\frac{r}{r^2}=H_s.
 \label{eq:weight-identities}
\end{equation}
Furthermore $w_t\le t^{-2}+\int_t^\infty x^{-2}dx\le2/t$, hence
$\sum_tw_t^2\le2\pi^2/3$.

Every pair is absent with probability $p=(1-1/n)^m$. For $n\ge m\ge1$, $p\ge1/4$ except when
$n=m=1$. For distinct pairs,
\[
 \E[I_iI_j]=(1-2/n)^m\le(1-1/n)^{2m}=p^2,
\]
so their covariance is nonpositive. Consequently
\[
 \E Z=q pH_s,
 \qquad
 \E Z^2\le(\E Z)^2+p q\sum_tw_t^2
 \le(\E Z)^2+\frac{2\pi^2pq}{3}.
\]
Since $q\ge1$, $p\ge1/4$, and $H_s\ge1$, Paley--Zygmund implies
\begin{equation}
 \Pp\{Z\ge \E Z/2\}
 \ge\frac{1}{4}\frac{(\E Z)^2}{\E Z^2}
 \ge c_{\rm occ}
 \label{eq:occupancy-pz}
\end{equation}
for a universal $c_{\rm occ}>0$. On this event,
\begin{equation}
 V\ge \frac{\eta^2(1-\eta)p}{16}\frac{H_s}{s}
 \ge c\frac{H_s}{s}.
 \label{eq:v-event}
\end{equation}

\subsection{A shifted Rademacher small-ball lemma}

\begin{lemma}\label{lem:shifted-smallball}
Let $\xi_1,\ldots,\xi_N$ be independent Rademachers, let $v_1,\ldots,v_N,b$ lie in a real
Hilbert space, and put $Y=\sum_i\xi_iv_i$ and $Q=\|Y-b\|^2$. Then
\[
 \E Q^2\le3(\E Q)^2,
 \qquad
 \Pp\{Q\ge\E Q/2\}\ge1/12.
\]
\end{lemma}

\begin{proof}
Write $v=\E\|Y\|^2=\sum_i\|v_i\|^2$ and $B=\|b\|^2$. The Hilbert-space fourth-moment
inequality gives $\E\|Y\|^4\le3v^2$. One direct proof expands the fourth power: terms with an
index of odd multiplicity vanish, while each surviving cross term is bounded by the corresponding
product of squared norms. Also
$\E\langle Y,b\rangle^2=\sum_i\langle v_i,b\rangle^2\le vB$.
Symmetry makes all cubic terms vanish. Expanding $Q^2$ therefore gives
\[
 \E Q^2
 =\E\|Y\|^4+B^2+4\E\langle Y,b\rangle^2+2Bv
 \le3v^2+B^2+6Bv
 \le3(v+B)^2.
\]
Since $\E Q=v+B$, the first claim follows. Paley--Zygmund applied to the nonnegative variable
$Q$ with threshold $1/2$ gives the second.
\end{proof}

Condition on sample pair indices, revealed signs, and learner randomness. The vector of all target
values at the first suffix vertices is an affine image of the missing independent signs. Define $Q$
by weighting the squared errors on those vertices uniformly by the lower bound
$(1-\eta)/(2n)$. The actual population error dominates $Q$, while
$Q=\|A\xi-b\|^2$ in a fixed weighted Euclidean space. Its conditional expectation is at least the
integrated posterior variance $V$. By
Lemma~\ref{lem:shifted-smallball}, $Q\ge V/2$ with conditional probability at least $1/12$.
Combining with \eqref{eq:occupancy-pz} and \eqref{eq:v-event} gives joint probability at least
$c_{\rm occ}/12$ that
\[
 L_{D_\sigma}(h_S)\ge c\frac{H_s}{s}.
\]
This probability is averaged over the uniform sign prior. Hence some deterministic $\sigma$ has
the same frequentist lower bound. Finally,
$H_s/s\ge c\min\{1,(q/m)\log(1+m/q)\}$ because $s=\lceil m/q\rceil$. This proves
Theorem~\ref{thm:directed-lower} except at $q=m=1$. At that endpoint, use the same two marginals
on the first pair without revealing its sign. Their one-sample total variation is $\eta<1$, while
their conditional averages differ by $\eta$. Le Cam's two-point argument gives a universal
constant-probability constant-risk lower bound, which is the required right side of
\eqref{eq:directed-lower} after changing constants.

\subsection{Inverting the risk lower bound}

For completeness, the condition
$m<cq\log(1/\epsilon)/\epsilon$ with a sufficiently small universal $c$ implies
$(q/m)\log(1+m/q)>c'\epsilon$ for all sufficiently small $\epsilon$. Bounded values of
$\epsilon$ are absorbed by changing constants. Thus Theorem~\ref{thm:directed-lower} implies
$\Omega(q\log(1/\epsilon)/\epsilon)$ constant-confidence sample complexity.

\section{Undirected lower bound}\label{app:undirected-lower}

We prove Theorem~\ref{thm:undirected-lower}. For the nontrivial regime $m\ge q$, fix
$\eta=a\sqrt{q/m}\le1/2$ for a sufficiently small universal $a$. Put the uniform prior on
$\sigma\in\{-1,+1\}^q$. Recall that the two endpoints of edge $i$ both have target
$\theta_i=(1-\eta\sigma_i)/2$.

For sign arrays differing only at coordinate $i$, the one-sample KL divergence is
\begin{align*}
 \KL(D_\sigma\|D_{\sigma^{(i)}})
 &=\frac{\eta}{q}\log\frac{1+\eta}{1-\eta}
 \le\frac{4\eta^2}{q}.
\end{align*}
Tensorization and Pinsker give
\[
 \TV(D_\sigma^m,D_{\sigma^{(i)}}^m)
 \le\sqrt{2m\eta^2/q}\le \sqrt2a.
\]
Choose $a\le1/(4\sqrt2)$. Assouad's argument then shows that for every decoder
$\widehat\sigma$,
\begin{equation}
 \E_{\sigma,S}d_H(\widehat\sigma,\sigma)\ge\frac{3q}{8}.
 \label{eq:assouad}
\end{equation}

Given a predictor $h$, decode $\widehat\sigma_i=+1$ if the average of $h(u_i)$ and $h(v_i)$ is
below $1/2$, and $-1$ otherwise. If this decoder errs, convexity gives
\[
 (h(u_i)-\theta_i)^2+(h(v_i)-\theta_i)^2\ge\frac{\eta^2}{2}.
\]
Every endpoint has mass at least $(1-\eta)/(2q)\ge1/(4q)$, so
\begin{equation}
 L_{D_\sigma}(h)\ge\frac{\eta^2}{8q}
 d_H(\widehat\sigma,\sigma).
 \label{eq:hamming-to-risk}
\end{equation}
Equations \eqref{eq:assouad}--\eqref{eq:hamming-to-risk} give Bayes expected risk at least
$3\eta^2/64=\Omega(q/m)$.

To obtain a probability statement, one may use a direct bounded-Hamming argument. Since
$0\le d_H\le q$ and its mean is at least $3q/8$,
\[
 \Pp\{d_H\ge3q/16\}\ge
 \frac{(3/8)-(3/16)}{1-(3/16)}=\frac3{13}.
\]
Together with \eqref{eq:hamming-to-risk}, the joint prior-and-sample probability of risk
$\Omega(q/m)$ is at least $3/13$. Averaging selects one deterministic $\sigma$. If $m<q$, set
$\eta$ to a small constant and apply the same proof to obtain constant risk, completing the stated
$\min\{1,q/m\}$ form.

\section{Auxiliary audits}

The supplementary program implements the following independent checks.
\begin{enumerate}[leftmargin=*,itemsep=2pt]
\item It enumerates every sample-location sequence for small $(q,s,m)$ and confirms
$\E V=pH_s/(4s)$ for the simpler uniform-label version of the chain calculation.
\item It verifies $\sum_tw_t=H_s$ and the uniform square-sum bound through $s=999$.
\item It enumerates all occupancy patterns induced by samples for small chains and checks the
constant-probability threshold in \eqref{eq:occupancy-pz}.
\item It enumerates Rademacher signs for random linear maps and shifts and checks
$\E Q^2\le3(\E Q)^2$.
\item It compares the weighted Caro--Wei sum with an exact exhaustive computation of the
independence number on random undirected graphs up to eight vertices and weights spanning six
orders of magnitude.
\item It exhausts the actual paired marginals and verifies their normalization, suffix targets,
posterior-variance identity, and the neighboring KL formula used in the edge lower bound.
\end{enumerate}
